\scanpage{11}
\margin{10}
\mspage{5}
est continue. Il revient au $=$ que l'application
\[
(15)\qquad\begin{aligned}
\operatorname{Aut}(Y)&\mathbin{\unsure{-|\longrightarrow}}\operatorname{Hom}(X,Z)\\
\varphi&\longmapsto v\varphi u
\end{aligned}
\]
soit continue.

Considérons \deleted{alors le} le stabilisateur \deleted{de}
$G_u$ de $u\in\operatorname{Hom}(X,Y)$ dans $\operatorname{Aut}(X)\times\operatorname{Aut}(Y)$
\[
\begin{aligned}
G_u&\subset\operatorname{Aut}(X)\times\operatorname{Aut}(Y)\\
\rotatebox{90}{$=$}\quad&\\[-1ex]
\{\psi,\varphi\mid\varphi u\psi^{-1}=u\quad\text{i.e.}\quad\varphi u=u\psi\}&
\end{aligned}
\]
c'est un sous-groupe fermé de $\operatorname{Aut}(X)\times\operatorname{Aut}(Y)$,
et l'application composée
\diagram{\begin{tikzcd}[ampersand replacement=\&]
G_u\arrow[r,"\mathrm{pr}_2"]\&\operatorname{Aut}(Y)\arrow[r,no head]\&\operatorname{Hom}(X,Z)
\end{tikzcd}}
n'est autre que $(\psi,\varphi)\longmapsto v\varphi u=\uline{vu}\psi$, qui est
donc continue, par construction de la top. de
$\operatorname{Hom}(X,Z)$. Donc l'application $\operatorname{Aut}(Y)\longrightarrow\operatorname{Hom}(X,Z)$
est continue en tous cas sur l'image \inserted{$H_u\subset\operatorname{Aut}(Y)$} de $G_u$
par $\mathrm{pr}_2$, \inserted{(car la top. de cette image est quotient de celle de $G_u$)} donc aussi sur tte réunion
finie de translatés à gauche \inserted{$\alpha_iH_u$ $(\alpha_i\in\operatorname{Aut}(Y))$} de cette image,
puisque l'on aura \deleted{\unsure{on}}
\[
v(\alpha_i\varphi)u=v\alpha_i u\psi.
\]
\margin{\uline{NB} $H_u$ est un sous-groupe \uuline{fermé} de $\operatorname{Aut}(Y)$, donc il est ouvert s'il est d'indice fini.}
Donc la continuité de (14), donc de (13), est
acquise, si on fait l'hypothèse :
\[
(16)\quad(*)\quad\left\{
\begin{minipage}{0.83\linewidth}
Pour toute flèche $u:X\mathbin{-}Y$ dans $C$, l'image
$H_u$ dans $\operatorname{Aut}_C(Y)$ du groupe $G_u$ des automs de cette
\inserted{flèche} est un sous-groupe d'\uline{indice fini de}
$\operatorname{Aut}(Y)$.
\end{minipage}\right.
\]
\margin{\uline{NB} L'hypothèse (16) implique que \inserted{dans} $\operatorname{Hom}(X,Y)$, les \deleted{\unsure{classes}} orbites \uline{sous} $\operatorname{Aut}(X)$ sont déjà ouvertes, \unsure{tte orbite sous|et tte orbite sous} $\operatorname{Aut}_{\unsure{0}}(X)\times\operatorname{Aut}_{\unsure{1}}(Y)$ \unsure{étant|est} réunion finie de telles orbites\ldots}
